\documentclass[11pt,oneside]{article}
\usepackage[a4paper,margin=27mm,headheight=15pt]{geometry}
\usepackage[T1]{fontenc}
\usepackage[utf8]{inputenc}
\IfFileExists{libertinus.sty}{\usepackage{libertinus}}{\usepackage{lmodern}}
\usepackage{microtype}
\usepackage{amsmath,amssymb,amsthm,mathtools,mathrsfs}
\usepackage{bm}
\usepackage{booktabs,longtable,array,tabularx,multirow}
\usepackage{graphicx}
\usepackage{xcolor}
\usepackage{enumitem}
\usepackage{fancyhdr}
\usepackage{titlesec}
\usepackage{listings}
\usepackage{xurl}
\usepackage{hyperref}
\usepackage[nameinlink,noabbrev]{cleveref}
\definecolor{linkblue}{RGB}{25,75,135}
\definecolor{shadegray}{RGB}{246,247,249}
\definecolor{rulegray}{RGB}{195,201,208}
\definecolor{newgreen}{RGB}{20,105,66}
\definecolor{importblue}{RGB}{36,79,128}
\definecolor{conjectureorange}{RGB}{151,83,15}
\hypersetup{colorlinks=true,linkcolor=linkblue,citecolor=linkblue,urlcolor=linkblue,
 pdftitle={Arithmetic Dyadic Rays of the Rvachev Fourier Product},
 pdfsubject={New cyclotomic decay laws for the Fabius--Rvachev--Thue--Morse system},
 pdfkeywords={Fabius function,Rvachev up-function,Thue-Morse,sinc product,cyclotomic norm,q-binomial}}
\pagestyle{fancy}
\fancyhf{}
\fancyhead[L]{\small Arithmetic dyadic rays of the Rvachev Fourier product}
\fancyhead[R]{\small\nouppercase{\leftmark}}
\fancyfoot[C]{\thepage}
\renewcommand{\headrulewidth}{0.4pt}
\titleformat{\section}{\Large\bfseries}{\thesection}{0.7em}{}
\titleformat{\subsection}{\large\bfseries}{\thesubsection}{0.7em}{}
\titleformat{\subsubsection}{\normalsize\bfseries}{\thesubsubsection}{0.7em}{}
\setcounter{secnumdepth}{3}
\setcounter{tocdepth}{2}
\setlist[itemize]{leftmargin=2em,itemsep=0.25em,topsep=0.4em}
\setlist[enumerate]{leftmargin=2.3em,itemsep=0.25em,topsep=0.4em}
\setlength{\emergencystretch}{3em}
\newtheorem{theorem}{Theorem}[section]
\newtheorem{proposition}[theorem]{Proposition}
\newtheorem{lemma}[theorem]{Lemma}
\newtheorem{corollary}[theorem]{Corollary}
\newtheorem{conjecture}[theorem]{Conjecture}
\theoremstyle{definition}
\newtheorem{definition}[theorem]{Definition}
\newtheorem{algorithm}[theorem]{Algorithm}
\newtheorem{example}[theorem]{Example}
\theoremstyle{remark}
\newtheorem{remark}[theorem]{Remark}
\newtheorem{warning}[theorem]{Warning}
\newcommand{\R}{\mathbb R}
\newcommand{\C}{\mathbb C}
\newcommand{\Q}{\mathbb Q}
\newcommand{\N}{\mathbb N}
\newcommand{\Z}{\mathbb Z}
\newcommand{\Up}{\operatorname{up}}
\newcommand{\sinc}{\operatorname{sinc}}
\newcommand{\ord}{\operatorname{ord}}
\newcommand{\Gal}{\operatorname{Gal}}
\newcommand{\Tr}{\operatorname{Tr}}
\newcommand{\Res}{\operatorname{Res}}
\newcommand{\Norm}{\operatorname{N}}
\newcommand{\e}{\mathrm e}
\newcommand{\ii}{\mathrm i}
\newcommand{\dd}{\,\mathrm d}
\newcommand{\bigO}{\mathcal O}
\newcommand{\smallo}{o}
\newcommand{\eps}{\varepsilon}
\newcommand{\wt}{s_2}
\newcommand{\qbinom}[3]{\genfrac{[}{]}{0pt}{}{#1}{#2}_{#3}}
\newcommand{\newresult}{\textcolor{newgreen}{\textsc{New result proved here.}}}
\newcommand{\imported}{\textcolor{importblue}{\textsc{Imported result.}}}
\newcommand{\openstatus}{\textcolor{conjectureorange}{\textsc{Open or conjectural.}}}
\newenvironment{boxedremark}{\par\smallskip\noindent\begin{center}\begin{minipage}{0.94\textwidth}\hrule\vspace{0.6em}\small}{\vspace{0.6em}\hrule\end{minipage}\end{center}\smallskip}
\lstdefinestyle{pseudo}{basicstyle=\ttfamily\footnotesize,backgroundcolor=\color{shadegray},frame=single,
 rulecolor=\color{rulegray},columns=fullflexible,keepspaces=true,showstringspaces=false,breaklines=true,
 xleftmargin=0.7em,xrightmargin=0.7em,aboveskip=0.8em,belowskip=0.8em}
\title{\bfseries Arithmetic Dyadic Rays of the Rvachev Fourier Product\\[0.35em]
\large Cyclotomic orbit norms, exact periodic decay, and a half-base $q$-binomial bridge}
\author{Frontier research report based on the complete \texttt{Analysis/FabiusFunction/docs} corpus\\
\small Vladimir Reshetnikov's \texttt{ProveIt} repository}
\date{27 August 2026}

\begin{document}
\maketitle

\begin{abstract}
The documentation corpus under \path{Analysis/FabiusFunction/docs} contains a mature,
Lean-backed exposition of the Fabius function and Rvachev's compactly supported
up-function, a large quarantine volume of non-formalized results, a Thue--Morse
formula atlas, twelve competing and synthesized studies of the Fourier decay of
$\Up$, seven source-paper transcriptions, and thirty-two frozen historical TeX
revisions. This report audits all fifty-seven TeX corpus items, distinguishes
results already present from genuinely unfilled directions, and develops an
arithmetic theory for
\[
 \mathscr U(x):=\widehat{\Up}(x)
 =\prod_{n=0}^{\infty}\sinc\!\left(\frac{\pi x}{2^n}\right).
\]
The central result is an exact decay law on every non-dyadic rational ray
$x_N=2^N(1+a/q)$. If $q$ is odd, $d=\ord_q(2)$, and
\[
 A_{a,q}=\left|\prod_{j=0}^{d-1}(1-\zeta_q^{a2^j})\right|,
\]
then
\[
 \log|\mathscr U(x_N)|
 =-\frac{(\log x_N)^2}{2\log2}-\kappa_{a,q}\log x_N+\Omega_{a,q}(N),
 \qquad
 \kappa_{a,q}=\frac12+\frac{\log(2\pi)}{\log2}-\frac{\log A_{a,q}}{d\log2},
\]
where $\Omega_{a,q}$ is explicitly periodic in $N$. The multiplier is a relative
cyclotomic norm of a Thue--Morse polynomial value. We construct its integral orbit
polynomial by a resultant, express its power traces through Thue--Morse signs and
Ramanujan sums, and prove an exact mean-exponent law over Galois orbits. A short
periodic-orbit argument gives the sharp bound $A_{a,q}^{1/d}\le\sqrt3$, with equality
only for the orbit $\{1/3,2/3\}$.

For the denominators $2^d-1$ and $2^d+1$ we obtain a refined two-sided expansion.
Both exponents grow as $d/2+1-(\log h)/(d\log2)$, where
$h=\prod_{k\ge1}\sinc(\pi/2^k)=-\mathscr U'(1)$, while the first exponentially small
corrections have opposite signs. Combining this with
$\mathfrak M_d=\prod_{j=1}^d(2^j-1)=2^{d(d+1)/2}(1/2;1/2)_d$ yields a direct limit
identity joining the half-base $q$-Pochhammer normalization, a cyclotomic orbit
norm, and the Rvachev sinc tail. All statements marked ``new'' are proved below.
Novelty is claimed relative to the reviewed repository corpus and the targeted
literature search; the root-of-unity recurrence itself is explicitly credited to
Sobolewski's broader automatic-polynomial theory.
\end{abstract}

\begin{boxedremark}
\textbf{Proof-status convention.}
\imported\ means that a result is already in the repository or cited literature.
\newresult\ means that a complete proof is supplied here and the statement was not
found in the reviewed corpus or targeted literature search. \openstatus\ labels
questions and numerical conjectures. This convention prevents known root-of-unity
identities from being relabeled as discoveries merely because they acquire a new
application to the Rvachev product.
\end{boxedremark}

\tableofcontents
\newpage

\section{Executive summary and frontier boundary}\label{sec:executive}

\subsection{The contribution map}

\begin{table}[ht]
\centering\small
\begin{tabularx}{\textwidth}{@{}p{0.20\textwidth}X p{0.18\textwidth}@{}}
\toprule
Result & Content & Status \\
\midrule
Rational-ray factorization & Splits $\mathscr U(2^Ny)$ into a universal log-Gaussian factor, a smooth mantissa tail, and a finite Thue--Morse product. & \newresult \\
Exact periodic ray law & Gives $\kappa_{a,q}$ and an explicit period-$d$ correction $\Omega_{a,q}((N+1)\bmod d)$, with no error term. & \newresult \\
Cyclotomic orbit polynomial & Constructs $\mathscr P_{q,2}(X)\in\Z[X]$ and proves a resultant identity whose right side is a perfect $d$th power. & \newresult \\
Ramanujan trace bridge & Expresses all power traces of the orbit multipliers as finite Thue--Morse--Ramanujan sums. & \newresult \\
Mean exponent law & Computes the Galois-orbit mean of $\kappa_{a,q}$ from the single integer $\Phi_q(1)$. & \newresult \\
Sharp endpoint & Proves $A_{a,q}^{1/d}\le\sqrt3$, uniquely extremized by $q=3$, and obtains the exact minimum rational-ray exponent. & \newresult \\
Mersenne--Fermat twin expansion & Gives matching main terms and opposite first corrections for denominators $2^d\mp1$. & \newresult \\
$q$-binomial bridge constant & Shows that two Mersenne-normalized cyclotomic norms converge to $(1/2;1/2)_\infty h$. & \newresult \\
Root-of-unity multiplier & The recurrence $T_{n+d}(\zeta)=\Gamma T_n(\zeta)$ and fixed-field interpretation. & \imported\ \cite{sobolewski2019} \\
Metric cocycle theory & Exact variance, central limit theorem, law of the iterated logarithm, transfer operator, and multifractal discussion. & \imported\ from the Fourier-decay drafts \\
\bottomrule
\end{tabularx}
\caption{The result boundary used throughout the report.}\label{tab:main-results}
\end{table}

\subsection{Why arithmetic rays are the genuine gap}

The Fourier-decay subcorpus already goes far beyond the elementary statement that
$\widehat{\Up}$ decreases faster than every power. Across its twelve drafts it
develops the logarithmic cocycle for the doubling map, almost-everywhere behavior,
exact variance, central-limit and iterated-logarithm fluctuations, extremal
envelopes, transfer operators, and a moment spectrum. Re-deriving that material
would add polish rather than frontier mathematics. The Thue--Morse atlas likewise
contains
\[
 T_n(z)=\sum_{m=0}^{2^n-1}(-1)^{\wt(m)}z^m
       =\prod_{j=0}^{n-1}(1-z^{2^j}),
\]
its Mahler recursion, rarefied sums, roots-of-unity methods, cyclotomic factors,
and extensive arithmetic subsequence material.

The missing composition was
\[
 \text{Rvachev sinc product}
 \longleftrightarrow
 \text{doubling-orbit sine product}
 \longleftrightarrow
 \text{cyclotomic Thue--Morse norm}.
\]
For irrational mantissas the middle term is a nontrivial ergodic cocycle. For a
rational mantissa it is periodic, and periodicity promotes an asymptotic heuristic
to an exact finite-state law. That promotion is the organizing idea of this report.

\subsection{Corpus and novelty audit}

The review treated every TeX path listed in Appendix~\ref{app:inventory} as a corpus
item. The live primary exposition, frontier volume, formula atlas, Fourier-decay
studies, and source-paper transcriptions were reviewed at source level. The
thirty-two frozen archive items were compared as revision chains, so historical
duplicates were not counted as independent mathematical claims. The archive's
provenance notes were used to identify material unique to older revisions.

The external search concentrated on Fabius/Rvachev Fourier products, lacunary
trigonometric products, Thue--Morse polynomials at roots of unity, cyclotomic
properties of automatic polynomials, and exact dyadic arithmetic. Sobolewski's
root-of-unity recurrence is the closest precursor: it supplies the multiplier and
fixed-field structure, but the located work does not state the Rvachev rational-ray
law, the orbit resultant in the form used below, the exponent mean law, the sharp
$\sqrt3$ endpoint, or the Mersenne--Fermat/$q$-Pochhammer bridge. The novelty claim
is therefore deliberately limited: these deductions are new relative to the
reviewed corpus and were not found explicitly in the located literature.

\section{Established platform from the corpus}\label{sec:platform}

\subsection{Fabius, Rvachev, and the Fourier product}

The primary exposition carefully separates the bounded Fabius distribution
function $F$, the even compactly supported Rvachev function $\Up$, and the signed
global continuation governed by Thue--Morse signs \cite{proveit,ariascompact,ariasarith}.
For this report only $\Up$ and its Fourier transform are needed. In the repository
normalization,
\begin{equation}
 \mathscr U(x):=\widehat{\Up}(x)
 =\prod_{n=0}^{\infty}\sinc\!\left(\frac{\pi x}{2^n}\right),
 \qquad \sinc u=\frac{\sin u}{u},\quad\sinc0=1.
 \label{eq:U-product}
\end{equation}
The product is entire and obeys
\begin{equation}
 \mathscr U(x)=\sinc(\pi x)\mathscr U(x/2).
 \label{eq:U-refine}
\end{equation}
Repeated use of the double-angle identity also yields
\begin{equation}
 \mathscr U(x)=\prod_{m=1}^{\infty}\left(\cos\frac{\pi x}{2^m}\right)^m.
 \label{eq:U-cos}
\end{equation}

\subsection{The finite Thue--Morse polynomial}

Write $\eps_m=(-1)^{\wt(m)}$, where $\wt(m)$ counts the binary ones in $m$. The
block polynomial is
\begin{equation}
 T_n(z):=\sum_{m=0}^{2^n-1}\eps_mz^m
       =\prod_{j=0}^{n-1}(1-z^{2^j}).
 \label{eq:TM-poly}
\end{equation}
The factorization follows from independent binary digit choices and encodes Prouhet
cancellation: $z=1$ is a zero of order $n$, hence
\[
 \sum_{m<2^n}\eps_mm^r=0\qquad(0\le r<n).
\]
The bridge to the Fourier product is
\begin{equation}
 |T_n(\e^{2\pi\ii t})|
 =\prod_{j=0}^{n-1}|1-\e^{2\pi\ii2^jt}|
 =\prod_{j=0}^{n-1}|2\sin(\pi2^jt)|.
 \label{eq:T-Q}
\end{equation}

\subsection{The half-base $q$-calculus layer}

The exact dyadic-value theory uses the base $q=1/2$ Pochhammer symbol and Gaussian
binomial coefficients. In the repository normalization,
\begin{equation}
 (z;1/2)_n=\sum_{k=0}^n(-1)^k2^{-\binom{k}{2}}
 \qbinom{n}{k}{1/2}z^k.
 \label{eq:qbinom-expansion}
\end{equation}
With the Mersenne factorial
\begin{equation}
 \mathfrak M_n:=\prod_{j=1}^n(2^j-1),
 \label{eq:Mersenne-factorial}
\end{equation}
one has
\begin{equation}
 \qbinom{n}{k}{1/2}
 =\frac{\mathfrak M_n}{\mathfrak M_k\mathfrak M_{n-k}2^{k(n-k)}},
 \qquad
 (1/2;1/2)_n=2^{-n(n+1)/2}\mathfrak M_n.
 \label{eq:Mersenne-qbinom}
\end{equation}
The first $q$ here is an analytic base fixed at $1/2$; later $q$ denotes an odd
cyclotomic denominator. Section~\ref{sec:qbridge} gives a limit in which these two
roles meet naturally rather than notationally.

\subsection{What is imported and what is not}

The twelve Fourier-decay manuscripts repeatedly derive a decomposition of
$\log|\mathscr U(2^Ny)|$ into a universal quadratic term and a Birkhoff sum for
$x\mapsto2x\pmod1$. Their final syntheses include the typical linear exponent,
exact covariance calculations, a central limit theorem, a law of the iterated
logarithm, transfer-operator formulas, and a multifractal moment spectrum. Those
results are context, not new contributions here.

For a primitive $q$th root $\zeta_q$ and $d=\ord_q(2)$, the sequence
$T_n(\zeta_q^a)$ obeys a two-term recurrence of span $d$. This is an elementary
special case of Sobolewski's general recurrence theory for automatic-sequence
polynomials at roots of unity \cite{sobolewski2019}. The same work places the
multiplier in the fixed field generated by $\langle2\rangle$ and proves integrality
properties. We therefore mark the recurrence itself as imported. What follows from
inserting it into the exact Rvachev factorization is the new contribution.

\section{Exact fixed-mantissa factorization}\label{sec:factorization}

Fix
\[
 1<y<2,\qquad t:=y-1\in(0,1),\qquad x_N:=2^Ny,
\]
and define
\begin{equation}
 H(y):=\prod_{r=1}^{\infty}\sinc\!\left(\frac{\pi y}{2^r}\right),
 \qquad
 Q_n(t):=\prod_{j=0}^{n-1}|2\sin(\pi2^jt)|,
 \qquad Q_0(t):=1.
 \label{eq:H-Q-def}
\end{equation}
The tail $H$ is positive and real analytic on $(0,2)$. By \eqref{eq:T-Q},
\begin{equation}
 Q_n(t)=|T_n(\e^{2\pi\ii t})|.
 \label{eq:Q-as-T}
\end{equation}
Thus all ray-dependent arithmetic is concentrated in a finite Thue--Morse product.

\begin{theorem}[Exact ray factorization]\label{thm:exact-factorization}
For every $N\ge0$ and $y\in(1,2)$,
\begin{equation}
 \boxed{
 |\mathscr U(2^Ny)|
 =2^{-N(N+1)/2}(\pi y)^{-(N+1)}H(y)\,2^{-(N+1)}Q_{N+1}(y-1).}
 \label{eq:exact-factorization}
\end{equation}
The identity remains valid when both sides vanish.

\smallskip\noindent\newresult
\end{theorem}

\begin{proof}
Split the defining product at $n=N$ and put $k=N-n$ in the finite part:
\[
 \mathscr U(2^Ny)=\prod_{k=0}^{N}\sinc(\pi2^ky)
 \prod_{r=1}^{\infty}\sinc\!\left(\frac{\pi y}{2^r}\right).
\]
The second product is $H(y)$. Since $y=1+t$ and $2^k$ is an integer,
$|\sin(\pi2^ky)|=|\sin(\pi2^kt)|$. The denominator of the finite product is
\[
 \prod_{k=0}^{N}\pi2^ky=(\pi y)^{N+1}2^{N(N+1)/2},
\]
whereas its numerator is $2^{-(N+1)}Q_{N+1}(t)$. Multiplication proves the claim.
\end{proof}

Put
\begin{equation}
 \lambda:=\log2,
 \qquad L_N:=\log x_N=N\lambda+\log y,
 \qquad
 \kappa_0:=\frac12+\frac{\log(2\pi)}{\lambda}
 =3.1514961294723187980\ldots,
 \label{eq:lambda-kappa0}
\end{equation}
and define
\begin{equation}
 \Psi(y):=\log H(y)-\log(2\pi)-\frac12\log y
 +\frac{\log(2\pi)\log y}{\lambda}
 +\frac{(\log y)^2}{2\lambda}.
 \label{eq:Psi-def}
\end{equation}

\begin{corollary}[Master cocycle identity]\label{cor:master-cocycle}
If $Q_{N+1}(t)\ne0$, then
\begin{equation}
 \boxed{
 \log|\mathscr U(2^Ny)|
 =-\frac{L_N^2}{2\lambda}-\kappa_0L_N
 +\Psi(y)+\log Q_{N+1}(t).}
 \label{eq:master-cocycle}
\end{equation}
If $Q_{N+1}(t)=0$, both sides are interpreted as $-\infty$.

\smallskip\noindent\newresult
\end{corollary}

\begin{proof}
Take logarithms in \eqref{eq:exact-factorization}. The deterministic finite part is
\[
 -\frac{N(N+1)}2\lambda-(N+1)\log(2\pi y)+\log H(y).
\]
Substituting $N=(L_N-\log y)/\lambda$ and collecting powers of $L_N$ gives
\eqref{eq:master-cocycle}; the constant is exactly \eqref{eq:Psi-def}.
\end{proof}

\begin{remark}[Relation with the metric theory]
For variable $t$,
\[
 \log Q_n(t)=\sum_{j=0}^{n-1}\log|2\sin(\pi2^jt)|
\]
is the Birkhoff sum studied in the Fourier-decay drafts. The present report fixes
$t$ rational. The doubling orbit then closes, turning the Birkhoff sum into a
finite algebraic cycle.
\end{remark}

\section{Exact arithmetic laws on rational dyadic rays}\label{sec:rational-rays}

\subsection{The cycle multiplier}

Let $q>1$ be odd, $(a,q)=1$, and write
\begin{equation}
 d:=\ord_q(2),\qquad
 \zeta_q:=\e^{2\pi\ii/q},\qquad
 H_q:=\langle2\rangle\subseteq(\Z/q\Z)^\times.
 \label{eq:q-d-H}
\end{equation}
Define
\begin{equation}
 \Gamma_{a,q}:=\prod_{j=0}^{d-1}(1-\zeta_q^{a2^j})=T_d(\zeta_q^a),
 \qquad A_{a,q}:=|\Gamma_{a,q}|>0,
 \qquad \alpha_{a,q}:=\frac1d\log A_{a,q}.
 \label{eq:Gamma-A-alpha}
\end{equation}
The value depends only on the coset $aH_q$.

\begin{proposition}[Root-of-unity block recurrence]\label{prop:root-recurrence}
For every $n\ge0$,
\begin{equation}
 T_{n+d}(\zeta_q^a)=\Gamma_{a,q}T_n(\zeta_q^a).
 \label{eq:root-recurrence}
\end{equation}
Consequently, if $n=md+r$, $0\le r<d$, then
\begin{equation}
 Q_n(a/q)=A_{a,q}^mQ_r(a/q).
 \label{eq:Q-periodic-block}
\end{equation}

\smallskip\noindent\imported\ in its general automatic-polynomial setting; the
special-case proof is included.
\end{proposition}

\begin{proof}
Using \eqref{eq:TM-poly} and $2^d\equiv1\pmod q$,
\begin{align*}
 T_{n+d}(\zeta_q^a)
 &=\prod_{j=0}^{n+d-1}(1-\zeta_q^{a2^j})\\
 &=\left(\prod_{j=0}^{d-1}(1-\zeta_q^{a2^j})\right)
   \left(\prod_{j=d}^{n+d-1}(1-\zeta_q^{a2^j})\right)\\
 &=\Gamma_{a,q}\prod_{j=0}^{n-1}(1-\zeta_q^{a2^j}).
\end{align*}
Iteration and absolute values give \eqref{eq:Q-periodic-block}.
\end{proof}

\subsection{The exact log-quadratic law}

Set
\begin{equation}
 y_{a,q}:=1+\frac aq\in(1,2),\qquad x_N:=2^Ny_{a,q},
 \label{eq:rational-ray}
\end{equation}
and define the arithmetic ray exponent
\begin{equation}
 \boxed{
 \kappa_{a,q}:=\kappa_0-\frac{\alpha_{a,q}}{\lambda}
 =\frac12+\frac{\log(2\pi)}{\log2}
 -\frac{\log A_{a,q}}{d\log2}.}
 \label{eq:kappa-aq}
\end{equation}
For $0\le r<d$, put
\begin{equation}
 \Omega_{a,q}(r):=\Psi(y_{a,q})
 +\left(1-\frac{\log y_{a,q}}{\lambda}\right)\alpha_{a,q}
 +\log Q_r(a/q)-r\alpha_{a,q}.
 \label{eq:Omega-def}
\end{equation}

\begin{theorem}[Exact rational-ray decay law]\label{thm:exact-rational-ray}
For every $N\ge0$, let $r_N$ be the residue of $N+1$ modulo $d$ in
$\{0,\ldots,d-1\}$. Then
\begin{equation}
 \boxed{
 \log|\mathscr U(2^Ny_{a,q})|
 =-\frac{(\log x_N)^2}{2\log2}
 -\kappa_{a,q}\log x_N+\Omega_{a,q}(r_N).}
 \label{eq:exact-ray-law}
\end{equation}
The complete remainder is periodic in $N$ with period dividing $d$; there is no
asymptotic error term.

\smallskip\noindent\newresult
\end{theorem}

\begin{proof}
Write $n=N+1=md+r_N$. Proposition~\ref{prop:root-recurrence} gives
\[
 \log Q_{N+1}(a/q)
 =(N+1)\alpha_{a,q}+
 \log Q_{r_N}(a/q)-r_N\alpha_{a,q}.
\]
Since $N+1=(L_N-\log y_{a,q})/\lambda+1$, substitution into
\eqref{eq:master-cocycle} changes the coefficient of $L_N$ from $-\kappa_0$ to
$-\kappa_{a,q}$ and leaves \eqref{eq:Omega-def}.
\end{proof}

\begin{corollary}[Exact $d$-step recurrence along a ray]\label{cor:d-step-ray}
Under the hypotheses above,
\begin{equation}
 \boxed{
 |\mathscr U(2^{N+d}y_{a,q})|
 =A_{a,q}\,2^{-d(d+3)/2}\pi^{-d}(2^Ny_{a,q})^{-d}
 |\mathscr U(2^Ny_{a,q})|.}
 \label{eq:d-step-ray}
\end{equation}

\smallskip\noindent\newresult
\end{corollary}

\begin{proof}
Take the quotient of \eqref{eq:exact-factorization} at $N+d$ and $N$. The tail
cancels and $Q_{N+d+1}/Q_{N+1}=A_{a,q}$; simplifying the remaining powers gives
the stated factor.
\end{proof}

\begin{remark}[Discrete scale invariance]
After removal of the universal log-Gaussian and linear factors, the values on a
rational ray lie on a finite cycle. The scale is $2^d$, and the multiplier is not a
fitted phenomenological constant but the cyclotomic norm \eqref{eq:Gamma-A-alpha}.
\end{remark}

\subsection{All rational mantissas}

Every reduced rational $t\in(0,1)$ has a unique representation
\begin{equation}
 t=\frac{a}{2^eq},\qquad e\ge0,\quad q\text{ odd},\quad(a,2q)=1.
 \label{eq:even-denominator}
\end{equation}
For $n>e$,
\begin{equation}
 Q_n(t)=Q_e(t)Q_{n-e}(a/q).
 \label{eq:preperiod-Q}
\end{equation}

\begin{corollary}[Zero-or-periodic dichotomy]\label{cor:all-rational-rays}
Let $t$ be as in \eqref{eq:even-denominator}.
\begin{enumerate}
 \item If $q=1$, then $Q_n(t)=0$ for every $n>e$, and
 $\mathscr U(2^N(1+t))=0$ for every $N\ge e$.
 \item If $q>1$, the finite prefix $Q_e(t)$ is followed by the exact law of
 \cref{thm:exact-rational-ray}, with period $\ord_q(2)$.
\end{enumerate}
Thus a rational ray is either eventually an exact zero ray or an eventually
periodic nonzero ray.

\smallskip\noindent\newresult
\end{corollary}

\begin{proof}
Split the product defining $Q_n$ at $j=e$. If $q=1$, the factor at $j=e$ is
$2|\sin(\pi a)|=0$. If $q>1$, the subsequent reduced odd-denominator orbit never
hits zero and \eqref{eq:Q-periodic-block} applies.
\end{proof}

\begin{example}[The denominator-three ray]
For $q=3$, $d=2$ and
\[
 \Gamma_{1,3}=(1-\zeta_3)(1-\zeta_3^2)=3,
 \qquad A_{1,3}=3.
\]
Hence
\[
 \kappa_{1,3}=\frac12+\frac{\log(2\pi/\sqrt3)}{\log2}
 =2.3590148791117407073\ldots.
\]
The periodic correction has at most two values.
\end{example}

\begin{example}[A neutral cycle at $q=15$]
The subgroup $\langle2\rangle$ has order $4$. Both cosets have multiplier $-1$,
so $A_{a,15}=1$ and the linear exponent is exactly $\kappa_0$. The orbit polynomial
is $(X+1)^2$, showing that it can have repeated roots when the multiplier lies in
a smaller subfield.
\end{example}

\begin{example}[The first visible split at $q=17$]
Here $d=8$ and the multiplier polynomial is
\[
 X^2-34X+17.
\]
Its roots are $17\pm4\sqrt{17}$, with amplitudes
$0.50757749752936\ldots$ and $33.492422502471\ldots$. The corresponding ray
exponents are $3.2737836274\ldots$ and $2.5182757764\ldots$.
\end{example}

\section{Cyclotomic algebra of the ray spectrum}\label{sec:cyclotomic-algebra}

\subsection{Relative norms and orbit polynomials}

Let $K_q=\Q(\zeta_q)$, $G_q=(\Z/q\Z)^\times$, and let $K_q^{H_q}$ be the fixed
field of $H_q=\langle2\rangle$. Then
\begin{equation}
 \Gamma_{a,q}=\Norm_{K_q/K_q^{H_q}}(1-\zeta_q^a).
 \label{eq:relative-norm}
\end{equation}
This imported fixed-field interpretation explains that every multiplier is an
algebraic integer. The quotient $G_q/H_q$ indexes its conjugates, with possible
repetitions. Write
\begin{equation}
 c_q^*:=|G_q/H_q|=\frac{\varphi(q)}d.
 \label{eq:coset-count}
\end{equation}

\begin{definition}[Cyclotomic orbit polynomial]
Choose a representative $a_C$ from every coset $C\in G_q/H_q$ and set
\begin{equation}
 \mathscr P_{q,2}(X):=\prod_{C\in G_q/H_q}(X-\Gamma_{a_C,q}).
 \label{eq:orbit-polynomial}
\end{equation}
\end{definition}

\begin{theorem}[Integrality and resultant construction]\label{thm:orbit-resultant}
Let $F_d(z):=T_d(z)=\prod_{j=0}^{d-1}(1-z^{2^j})$. Then
\begin{equation}
 \mathscr P_{q,2}(X)\in\Z[X]
 \label{eq:orbit-integral}
\end{equation}
and
\begin{equation}
 \boxed{
 \Res_z\!\left(\Phi_q(z),X-F_d(z)\right)=\mathscr P_{q,2}(X)^d.}
 \label{eq:resultant-power}
\end{equation}
The polynomial on the right need not be squarefree.

\smallskip\noindent\newresult
\end{theorem}

\begin{proof}
The coefficients of \eqref{eq:orbit-polynomial} are symmetric polynomials in
algebraic integers, hence algebraic integers. The Galois group permutes the cosets,
so every coefficient lies in $\Q$ and therefore in $\Z$.

Since $\Phi_q$ is monic with roots $\zeta_q^a$, $a\in G_q$,
\[
 \Res_z(\Phi_q(z),X-F_d(z))
 =\prod_{a\in G_q}(X-F_d(\zeta_q^a)).
\]
The value $F_d(\zeta_q^a)=\Gamma_{a,q}$ is constant on a coset $aH_q$, and every
coset contains $d$ elements. Grouping factors proves \eqref{eq:resultant-power}.
\end{proof}

\subsection{Ramanujan traces}

Recall the Ramanujan sum
\begin{equation}
 c_q(m):=\sum_{a\in G_q}\zeta_q^{am}
 =\mu\!\left(\frac{q}{(q,m)}\right)
 \frac{\varphi(q)}{\varphi(q/(q,m))}.
 \label{eq:Ramanujan-sum}
\end{equation}

\begin{theorem}[Power traces as finite digital sums]\label{thm:Ramanujan-traces}
For every integer $r\ge1$,
\begin{equation}
 \boxed{
 \sum_{C\in G_q/H_q}\Gamma_{a_C,q}^{\,r}
 =\frac1d\sum_{0\le m_1,\ldots,m_r<2^d}
 \eps_{m_1}\cdots\eps_{m_r}
 c_q(m_1+\cdots+m_r).}
 \label{eq:power-trace}
\end{equation}
In particular,
\begin{equation}
 \Tr_{K_q^{H_q}/\Q}(\Gamma_{1,q})
 =\frac1d\sum_{m=0}^{2^d-1}\eps_mc_q(m),
 \label{eq:first-trace}
\end{equation}
with the natural embedding multiplicities if $\Gamma_{1,q}$ lies in a smaller
subfield.

\smallskip\noindent\newresult
\end{theorem}

\begin{proof}
Expand
\[
 \Gamma_{a,q}=T_d(\zeta_q^a)=\sum_{m=0}^{2^d-1}\eps_m\zeta_q^{am}.
\]
The value is constant on each $H_q$-coset, so summing over cosets is $1/d$ times
summing over all $a\in G_q$. Raise the finite sum to the $r$th power, interchange
sums, and apply \eqref{eq:Ramanujan-sum}.
\end{proof}

\begin{corollary}[Newton reconstruction]\label{cor:Newton-reconstruction}
The first $c_q^*$ power sums in \eqref{eq:power-trace}, together with Newton's
identities, reconstruct $\mathscr P_{q,2}(X)$ using integer arithmetic only.
This is an independent check on the resultant method.
\end{corollary}

\subsection{A one-integer law for the mean exponent}

\begin{theorem}[Norm, constant term, and mean ray exponent]\label{thm:mean-exponent}
With $c_q^*=\varphi(q)/d$,
\begin{align}
 \prod_{C\in G_q/H_q}\Gamma_{a_C,q}&=\Phi_q(1),
 \label{eq:product-Gamma}\\
 \mathscr P_{q,2}(0)&=(-1)^{c_q^*}\Phi_q(1),
 \label{eq:P-constant}\\
 \prod_{C\in G_q/H_q}A_{a_C,q}&=\Phi_q(1),
 \label{eq:product-A}
\end{align}
and
\begin{equation}
 \boxed{
 \frac1{c_q^*}\sum_{C\in G_q/H_q}\kappa_{a_C,q}
 =\kappa_0-\frac{\log\Phi_q(1)}{\varphi(q)\log2}.}
 \label{eq:mean-kappa}
\end{equation}
For $q>1$,
\begin{equation}
 \Phi_q(1)=
 \begin{cases}
 p,&q=p^s\text{ is a prime power},\\
 1,&q\text{ has at least two distinct prime divisors}.
 \end{cases}
 \label{eq:Phi-one}
\end{equation}
Thus every non-prime-power denominator has mean exponent exactly $\kappa_0$.

\smallskip\noindent\newresult
\end{theorem}

\begin{proof}
By transitivity of norms,
\[
 \prod_C\Gamma_{a_C,q}
 =\Norm_{K_q^{H_q}/\Q}
 \!\left(\Norm_{K_q/K_q^{H_q}}(1-\zeta_q)\right)
 =\Norm_{K_q/\Q}(1-\zeta_q)=\Phi_q(1).
\]
This proves \eqref{eq:product-Gamma}; the constant term and absolute-value product
follow. Taking logarithms in \eqref{eq:product-A}, inserting \eqref{eq:kappa-aq},
and using $dc_q^*=\varphi(q)$ gives \eqref{eq:mean-kappa}. Formula
\eqref{eq:Phi-one} is the classical value of a cyclotomic polynomial at one.
\end{proof}

\begin{remark}[Compensation between conjugate rays]
For a denominator with at least two prime factors, a coset with $A_{a,q}>1$ and
hence a smaller-than-$\kappa_0$ exponent must be compensated geometrically by
cosets with $A_{a,q}<1$. The $q=17$ and $q=31$ tables below show this splitting
explicitly.
\end{remark}

\section{The sharp rational-ray endpoint}\label{sec:sharp-endpoint}

\begin{theorem}[Sharp periodic-orbit sine-product bound]\label{thm:sqrt3-bound}
For every odd $q>1$ and every $(a,q)=1$,
\begin{equation}
 \boxed{A_{a,q}^{1/d}\le\sqrt3.}
 \label{eq:sqrt3-bound}
\end{equation}
Equality holds if and only if $q=3$ (with $a=1$ or $2$).

\smallskip\noindent\newresult
\end{theorem}

\begin{proof}
Let
\[
 t_j=\left\{\frac{a2^j}{q}\right\},\qquad
 c_j=\cos(2\pi t_j),\qquad0\le j<d,
\]
with indices cyclic. Doubling gives $c_{j+1}=2c_j^2-1$. If
$s=d^{-1}\sum_jc_j$, cyclic averaging and Cauchy--Schwarz give
\[
 s=2\frac1d\sum_jc_j^2-1\ge2s^2-1.
\]
Hence $(2s+1)(s-1)\le0$ and $s\ge-1/2$. By arithmetic--geometric mean,
\begin{align*}
 A_{a,q}^{2/d}
 &=\left(\prod_{j=0}^{d-1}|1-\e^{2\pi\ii t_j}|^2\right)^{1/d}\\
 &=\left(\prod_{j=0}^{d-1}(2-2c_j)\right)^{1/d}
 \le\frac1d\sum_{j=0}^{d-1}(2-2c_j)
 =2-2s\le3.
\end{align*}
If equality holds, then $s=-1/2$ and equality in arithmetic--geometric mean forces
all $c_j=-1/2$. Every orbit point is therefore $1/3$ or $2/3$ modulo one, so the
reduced orbit is exactly $\{1/3,2/3\}$ and $q=3$. The converse is immediate.
\end{proof}

\begin{corollary}[Sharp minimum of the rational-ray exponent]\label{cor:kappa-minimum}
Every nonzero rational ray satisfies
\begin{equation}
 \boxed{
 \kappa_{a,q}\ge\kappa_{\min}:=\frac12+\frac{\log(2\pi/\sqrt3)}{\log2}
 =2.3590148791117407073\ldots.}
 \label{eq:kappa-minimum}
\end{equation}
Equality is attained only by the denominator-three doubling orbit.

\smallskip\noindent\newresult
\end{corollary}

\begin{proof}
Insert \eqref{eq:sqrt3-bound} into \eqref{eq:kappa-aq}.
\end{proof}

\begin{remark}[A spectral interpretation]
The arithmetic exponents form a spectrum of cycle averages of
$\log|2\sin(\pi t)|$. The theorem identifies the exact lower endpoint without
using transfer-operator machinery. Section~\ref{sec:mersenne} constructs exponents
tending to $+\infty$, so the spectrum has a sharp bottom and no finite top.
\end{remark}

\section{Mersenne and plus-denominator rays}\label{sec:mersenne}

The preceding theory turns a fast-decay construction into an arithmetic design
problem: choose a rational doubling cycle for which $A_{a,q}$ is very small. The
simplest such cycles come from the denominators immediately adjacent to a power of
two.

\subsection{Exact orbit products}

For $d\ge2$ put
\begin{equation}
 q_d^-:=2^d-1,
 \qquad q_d^+:=2^d+1,
 \label{eq:qpm}
\end{equation}
and let $A_d^-$ and $A_d^+$ denote the amplitudes $A_{1,q}$ for these denominators.

\begin{lemma}[Orders of two]\label{lem:orders-pm}
One has
\begin{equation}
 \ord_{q_d^-}(2)=d,
 \qquad
 \ord_{q_d^+}(2)=2d.
 \label{eq:orders-pm}
\end{equation}
\end{lemma}

\begin{proof}
For $q_d^-$, if $2^r\equiv1\pmod{2^d-1}$ with $0<r<d$, then the positive integer
$2^r-1$ would be divisible by the larger integer $2^d-1$, impossible.

For $q_d^+$, $2^d\equiv-1$, so the order $r$ is even. Inside the cyclic group
generated by $2$, the unique element of order two is $2^{r/2}$; hence
$d=(2m+1)r/2$ for some $m\ge0$. If $m\ge1$, then $r\le2d/3<d$, but
$q_d^+$ would divide $2^r-1<2^d+1$, again impossible. Thus $m=0$ and $r=2d$.
\end{proof}

\begin{proposition}[Exact product forms]\label{prop:pm-products}
For $d\ge2$,
\begin{align}
 A_d^-&=\prod_{k=1}^{d}
 2\sin\!\left(\frac{\pi2^{-k}}{1-2^{-d}}\right),
 \label{eq:Aminus-product}\\
 A_d^+&=\left[
 \prod_{k=1}^{d}
 2\sin\!\left(\frac{\pi2^{-k}}{1+2^{-d}}\right)
 \right]^2.
 \label{eq:Aplus-product}
\end{align}
\end{proposition}

\begin{proof}
For $q_d^-$, the orbit is $1,2,\ldots,2^{d-1}$ under multiplication by two.
Reindexing $j=d-k$ gives
\[
 \frac{2^j}{2^d-1}=\frac{2^{-k}}{1-2^{-d}}.
\]
For $q_d^+$ the order is $2d$ and $2^{j+d}\equiv-2^j\pmod{q_d^+}$, so the
second half of the orbit contributes the same absolute factors as the first.
The same reindexing proves \eqref{eq:Aplus-product}.
\end{proof}

\subsection{The universal sinc-tail constant}

Define
\begin{equation}
 B_d:=\prod_{k=1}^{d}2\sin\!\left(\frac{\pi}{2^k}\right),
 \qquad
 H_d(1):=\prod_{k=1}^{d}\sinc\!\left(\frac{\pi}{2^k}\right).
 \label{eq:B-Hd}
\end{equation}
Then
\begin{equation}
 B_d=(2\pi)^d2^{-d(d+1)/2}H_d(1).
 \label{eq:B-factor}
\end{equation}
The finite sinc tail converges to
\begin{equation}
 h:=\prod_{k=1}^{\infty}\sinc\!\left(\frac{\pi}{2^k}\right)
 =0.5537712758888100953442174163445488813\ldots.
 \label{eq:h-def}
\end{equation}

\begin{lemma}[Three forms of the tail constant]\label{lem:h-forms}
The constant $h$ satisfies
\begin{align}
 h&=-\mathscr U'(1),
 \label{eq:h-derivative}\\
 \log h&=-\sum_{m=1}^{\infty}
 \frac{\zeta(2m)}{m(4^m-1)}.
 \label{eq:h-zeta-series}
\end{align}
Moreover,
\begin{equation}
 \log H_d(1)=\log h+\bigO(4^{-d}).
 \label{eq:Hd-tail}
\end{equation}
\end{lemma}

\begin{proof}
From \eqref{eq:U-refine},
\[
 \mathscr U(x)=\sinc(\pi x)
 \prod_{k=1}^{\infty}\sinc\!\left(\frac{\pi x}{2^k}\right).
\]
The derivative of $\sinc(\pi x)$ at $x=1$ is $-1$, proving
\eqref{eq:h-derivative}. Euler's product for the sine gives, for $|z|<\pi$,
\[
 \log\sinc z=-\sum_{m=1}^{\infty}\frac{\zeta(2m)}{m\pi^{2m}}z^{2m}.
\]
Substitute $z=\pi/2^k$, sum over $k\ge1$, and use
$\sum_{k\ge1}4^{-km}=1/(4^m-1)$. Finally,
$\log\sinc(\pi/2^k)=\bigO(4^{-k})$, which yields \eqref{eq:Hd-tail}.
\end{proof}

The first exponentially small correction is controlled by the convergent constant
\begin{equation}
 C_{\cot}:=\sum_{k=1}^{\infty}
 \left(\pi2^{-k}\cot(\pi2^{-k})-1\right)
 =-1.2837124730461260445945438959726042474\ldots.
 \label{eq:Ccot}
\end{equation}
Convergence follows from $x\cot x=1-x^2/3+\bigO(x^4)$ at the origin.

\begin{theorem}[Refined twin product expansion]\label{thm:pm-product-asymptotic}
As $d\to\infty$,
\begin{align}
 \log\frac{A_d^-}{B_d}
 &=(d+C_{\cot})2^{-d}+\bigO(d4^{-d}),
 \label{eq:Aminus-refined}\\
 \log\frac{\sqrt{A_d^+}}{B_d}
 &=-(d+C_{\cot})2^{-d}+\bigO(d4^{-d}).
 \label{eq:Aplus-refined}
\end{align}
Equivalently,
\begin{align}
 A_d^-&=h(2\pi)^d2^{-d(d+1)/2}
 \exp\!\left((d+C_{\cot})2^{-d}+\bigO(d4^{-d})\right),
 \label{eq:Aminus-full}\\
 A_d^+&=h^2(2\pi)^{2d}2^{-d(d+1)}
 \exp\!\left(-2(d+C_{\cot})2^{-d}+\bigO(d4^{-d})\right).
 \label{eq:Aplus-full}
\end{align}

\smallskip\noindent\newresult
\end{theorem}

\begin{proof}
Put $\delta=2^{-d}$, $x_k=\pi2^{-k}$, and $g(x)=\log(2\sin x)$. From
\cref{prop:pm-products},
\begin{align*}
 \log(A_d^-/B_d)
 &=\sum_{k=1}^{d}\left[g\!\left(\frac{x_k}{1-\delta}\right)-g(x_k)\right],\\
 \log(\sqrt{A_d^+}/B_d)
 &=\sum_{k=1}^{d}\left[g\!\left(\frac{x_k}{1+\delta}\right)-g(x_k)\right].
\end{align*}
On the relevant interval, $xg'(x)=x\cot x$ and $x^2g''(x)$ are uniformly bounded.
Taylor expansion under a multiplicative perturbation therefore gives, uniformly in
$k$,
\[
 g\!\left(\frac{x_k}{1\mp\delta}\right)-g(x_k)
 =\pm\delta x_k\cot x_k+\bigO(\delta^2).
\]
Summing produces an error $\bigO(d\delta^2)$. Further,
\[
 \sum_{k=1}^{d}x_k\cot x_k
 =d+C_{\cot}+\bigO(4^{-d}).
\]
This proves \eqref{eq:Aminus-refined}--\eqref{eq:Aplus-refined}. Equations
\eqref{eq:Aminus-full}--\eqref{eq:Aplus-full} follow from
\eqref{eq:B-factor} and \eqref{eq:Hd-tail}.
\end{proof}

\subsection{Exponent asymptotics and unboundedness}

Let $\kappa_d^-$ and $\kappa_d^+$ be the ray exponents for $q_d^-$ and $q_d^+$,
respectively, and put
\begin{equation}
 K_d:=\frac d2+1-\frac{\log h}{d\log2}.
 \label{eq:Kd}
\end{equation}

\begin{corollary}[Opposite first corrections]\label{cor:pm-kappa}
As $d\to\infty$,
\begin{align}
 \boxed{
 \kappa_d^-=K_d-\frac{2^{-d}}{\log2}
 \left(1+\frac{C_{\cot}}d\right)+\bigO(4^{-d}),}
 \label{eq:kminus-asymptotic}\\
 \boxed{
 \kappa_d^+=K_d+\frac{2^{-d}}{\log2}
 \left(1+\frac{C_{\cot}}d\right)+\bigO(4^{-d}).}
 \label{eq:kplus-asymptotic}
\end{align}
In particular,
\begin{equation}
 \kappa_d^{\pm}=\frac d2+1-\frac{\log h}{d\log2}+\bigO(2^{-d}),
 \label{eq:kpm-main}
\end{equation}
and the rational-ray exponent spectrum is unbounded above.

\smallskip\noindent\newresult
\end{corollary}

\begin{proof}
Insert \eqref{eq:Aminus-full} into \eqref{eq:kappa-aq} with order $d$, and
\eqref{eq:Aplus-full} with order $2d$. The terms $\log(2\pi)$ cancel against
$\kappa_0$, and the remaining algebra gives \eqref{eq:kminus-asymptotic} and
\eqref{eq:kplus-asymptotic}. Since $K_d\sim d/2$, the exponents tend to infinity.
\end{proof}

Combining \cref{cor:kappa-minimum,cor:pm-kappa} gives a complete endpoint statement.

\begin{corollary}[Endpoint geometry of the rational spectrum]\label{cor:spectrum-endpoints}
Let $\mathcal K_{\mathrm{rat}}$ be the set of exponents of all nonzero rational
rays. Then
\begin{equation}
 \min\mathcal K_{\mathrm{rat}}=\kappa_{\min},
 \qquad
 \sup\mathcal K_{\mathrm{rat}}=+\infty.
 \label{eq:spectrum-endpoints}
\end{equation}
The minimum is attained only by denominator $3$.
\end{corollary}

\section{A direct half-base $q$-binomial bridge}\label{sec:qbridge}

The Mersenne factorial \eqref{eq:Mersenne-factorial} is the denominator-clearing
normalization behind the half-base Gaussian coefficients. The preceding product
asymptotics show that the same arithmetic factors normalize special cyclotomic
ray multipliers.

Set
\begin{equation}
 \mathcal C_*:=(1/2;1/2)_\infty h
 =0.1599225518776068325779014528557986690\ldots,
 \label{eq:Cstar}
\end{equation}
where
\begin{equation}
 (1/2;1/2)_\infty
 =0.2887880950866024212788997219292307801\ldots.
 \label{eq:qpoch-number}
\end{equation}

\begin{theorem}[Mersenne--cyclotomic--sinc bridge]\label{thm:qbridge}
As $d\to\infty$,
\begin{equation}
 \boxed{
 \frac{\mathfrak M_dA_d^-}{(2\pi)^d}\longrightarrow\mathcal C_*},
 \qquad
 \boxed{
 \frac{\mathfrak M_d\sqrt{A_d^+}}{(2\pi)^d}\longrightarrow\mathcal C_*}.
 \label{eq:qbridge-limits}
\end{equation}
More precisely,
\begin{align}
 \frac{\mathfrak M_dA_d^-}{(2\pi)^d}
 &=\mathcal C_*
 \exp\!\left((d+C_{\cot}+1)2^{-d}+\bigO(d4^{-d})\right),
 \label{eq:qbridge-minus-refined}\\
 \frac{\mathfrak M_d\sqrt{A_d^+}}{(2\pi)^d}
 &=\mathcal C_*
 \exp\!\left((1-d-C_{\cot})2^{-d}+\bigO(d4^{-d})\right).
 \label{eq:qbridge-plus-refined}
\end{align}

\smallskip\noindent\newresult
\end{theorem}

\begin{proof}
Using \eqref{eq:Mersenne-qbinom} and \eqref{eq:B-factor},
\[
 \frac{\mathfrak M_dB_d}{(2\pi)^d}
 =(1/2;1/2)_dH_d(1).
\]
Now multiply by $A_d^-/B_d$ or by $\sqrt{A_d^+}/B_d$ and use
\cref{thm:pm-product-asymptotic}. The limits follow because
$(1/2;1/2)_d\to(1/2;1/2)_\infty$ and $H_d(1)\to h$.

For the refined forms, note
\[
 \log\frac{(1/2;1/2)_d}{(1/2;1/2)_\infty}
 =-\sum_{k>d}\log(1-2^{-k})=2^{-d}+\bigO(4^{-d}),
\]
while \eqref{eq:Hd-tail} contributes only $\bigO(4^{-d})$. Adding the exponent
corrections from \eqref{eq:Aminus-refined} and \eqref{eq:Aplus-refined} proves
\eqref{eq:qbridge-minus-refined}--\eqref{eq:qbridge-plus-refined}.
\end{proof}

\begin{boxedremark}
\textbf{Interpretation.} The limit \eqref{eq:qbridge-limits} is a literal bridge
among the three recurrent structures of the repository:
\[
 \underbrace{\mathfrak M_d}_{\text{half-base Gaussian normalization}}
 \times
 \underbrace{A_d^-\text{ or }\sqrt{A_d^+}}_{\text{cyclotomic Thue--Morse orbit norm}}
 \times
 \underbrace{(2\pi)^{-d}}_{\text{sinc scaling}}
 \longrightarrow
 \underbrace{(1/2;1/2)_\infty(-\mathscr U'(1))}_{\text{global constant}}.
\]
This is more than a shared appearance of powers of two: it is a quantitative
asymptotic identity tying the $q$-binomial and Fourier sides together.
\end{boxedremark}

\subsection{Numerical convergence of the twin exponents}

\begin{table}[ht]
\centering\small
\begin{tabular}{rrrrrr}
\toprule
$d$ & $2^d-1$ & $\kappa_d^-$ & $2^d+1$ & $\kappa_d^+$ & $K_d$ \\
\midrule
2  & 3     & 2.35901487911 & 5     & 2.57101410575 & 2.42631893591 \\
3  & 7     & 2.68360364246 & 9     & 2.88733571269 & 2.78421262394 \\
4  & 15    & 3.15149612947 & 17    & 3.27378362740 & 3.21315946796 \\
5  & 31    & 3.63677701126 & 33    & 3.70378409640 & 3.67052757436 \\
6  & 63    & 4.12430649454 & 65    & 4.15974429455 & 4.14210631197 \\
7  & 127   & 4.61257768880 & 129   & 4.63098586366 & 4.62180541026 \\
8  & 255   & 5.10184195614 & 257   & 5.11130441973 & 5.10657973398 \\
9  & 511   & 5.59231992724 & 513   & 5.59715163452 & 5.59473754131 \\
10 & 1023  & 6.08403530020 & 1025  & 6.08649134438 & 6.08526378718 \\
11 & 2047  & 6.57689018061 & 2049  & 6.57813464441 & 6.57751253380 \\
12 & 4095  & 7.07073858328 & 4097  & 7.07136766593 & 7.07105315599 \\
13 & 8191  & 7.56542880068 & 8193  & 7.56574624039 & 7.56558752860 \\
14 & 16383 & 8.06082272205 & 16385 & 8.06098268408 & 8.06090270513 \\
\bottomrule
\end{tabular}
\caption{Mersenne and plus-denominator exponents bracketing the common main term
$K_d$. The gap has the predicted size $2^{1-d}/\log2$ to first order.}
\label{tab:pm-kappa}
\end{table}

\section{Exact computations and orbit-polynomial data}\label{sec:computations}

The theory is computationally concrete. For a given odd $q$ one needs only the
multiplicative order of $2$, a finite product at roots of unity, and either a
resultant or Ramanujan power sums. No numerical integration of $\Up$ is involved.

\subsection{Two exact algorithms}

\begin{algorithm}[Resultant route]\label{alg:resultant}
Given odd $q>1$:
\begin{enumerate}
 \item compute $d=\ord_q(2)$;
 \item form $F_d(z)=\prod_{j=0}^{d-1}(1-z^{2^j})$;
 \item compute $R(X)=\Res_z(\Phi_q(z),X-F_d(z))$ over $\Z[X]$;
 \item factor $R$ and divide every factor exponent by $d$;
 \item the resulting exact polynomial is $\mathscr P_{q,2}(X)$;
 \item take the absolute values of its complex roots and insert them into
 \eqref{eq:kappa-aq}.
\end{enumerate}
The perfect-$d$th-power assertion in \eqref{eq:resultant-power} is a built-in check.
\end{algorithm}

\begin{algorithm}[Ramanujan--Newton route]\label{alg:ramanujan}
Let $c_q^*=\varphi(q)/d$. Compute the power sums
$p_r=\sum_C\Gamma_C^r$ for $1\le r\le c_q^*$ by \eqref{eq:power-trace}, then apply
Newton's identities to recover the elementary symmetric functions and hence
$\mathscr P_{q,2}(X)$. This route uses integer arithmetic and can exploit fast
convolution of the finite Thue--Morse block.
\end{algorithm}

The resultant route is simplest for modest $d$. The Ramanujan route becomes
attractive when $2^d$ is too large to expand naively but $q$ has a simple divisor
structure, because $c_q(m)$ depends only on $(q,m)$.

\subsection{Exact orbit polynomials}

\begin{table}[ht]
\centering\small
\begin{tabularx}{\textwidth}{@{}r r X@{}}
\toprule
$q$ & $d=\ord_q(2)$ & $\mathscr P_{q,2}(X)$ \\
\midrule
3  & 2  & $X-3$ \\
5  & 4  & $X-5$ \\
7  & 3  & $X^2+7$ \\
9  & 6  & $X-3$ \\
11 & 10 & $X-11$ \\
15 & 4  & $(X+1)^2$ \\
17 & 8  & $X^2-34X+17$ \\
21 & 6  & $(X+1)^2$ \\
23 & 11 & $X^2+23$ \\
31 & 5  & $X^6+93X^4+899X^2+31$ \\
\bottomrule
\end{tabularx}
\caption{Exact orbit polynomials obtained by the resultant method. Repeated roots
at $q=15$ and $21$ reflect a larger stabilizer than $\langle2\rangle$.}
\label{tab:orbit-polynomials}
\end{table}

Several structural features are visible immediately.
\begin{itemize}
 \item If $2$ generates all of $G_q$, the orbit polynomial is linear and the
 multiplier is $\Phi_q(1)$; this explains $q=3,5,9,11$.
 \item For $q=7$ and $23$, the two multipliers are purely imaginary conjugates and
 both amplitudes equal $\sqrt q$.
 \item At $q=17$, the two multipliers are real and split as $17\pm4\sqrt{17}$.
 \item For $q=31$, the odd order $d=5$ makes the multiplier roots occur in imaginary
 pairs; the squared amplitudes are the positive roots of
 $Y^3-93Y^2+899Y-31$.
\end{itemize}

\begin{table}[ht]
\centering\footnotesize
\begin{tabularx}{\textwidth}{@{}r r X X@{}}
\toprule
$q$ & multiplicity & amplitude $A_{a,q}$ & exponent $\kappa_{a,q}$ \\
\midrule
3  & 1 & $3$ & $2.3590148791117$ \\
5  & 1 & $5$ & $2.5710141057505$ \\
7  & 2 & $\sqrt7=2.6457513110646\ldots$ & $2.6836036424627$ \\
9  & 1 & $3$ & $2.8873357126855$ \\
11 & 1 & $11$ & $2.8055529676086$ \\
15 & 2 & $1$ & $3.1514961294723$ \\
17 & 1 & $17-4\sqrt{17}=0.50757749752936\ldots$ & $3.2737836273969$ \\
17 & 1 & $17+4\sqrt{17}=33.492422502471\ldots$ & $2.5182757763915$ \\
21 & 2 & $1$ & $3.1514961294723$ \\
23 & 2 & $\sqrt{23}=4.7958315233127\ldots$ & $2.9458796769243$ \\
31 & 2 & $0.18602850302165\ldots$ & $3.6367770112610$ \\
31 & 2 & $3.3042170276289\ldots$ & $2.8066344396753$ \\
31 & 2 & $9.0580098934805\ldots$ & $2.5156573064420$ \\
\bottomrule
\end{tabularx}
\caption{Amplitude classes and rational-ray exponents. Multiplicity is the number
of $H_q$-cosets with that amplitude.}
\label{tab:ray-exponents}
\end{table}

\subsection{Exact checks performed}

The calculations behind \cref{tab:orbit-polynomials,tab:ray-exponents} were checked
in three independent ways:
\begin{enumerate}
 \item the symbolic resultant was factored and verified to be a perfect $d$th
 power;
 \item high-precision direct products $\prod_{j<d}(1-\zeta_q^{a2^j})$ were evaluated
 for every coset and substituted into the exact integer polynomial;
 \item the product of all orbit amplitudes was checked against $\Phi_q(1)$, and the
 numerical mean of the exponents was checked against \eqref{eq:mean-kappa}.
\end{enumerate}
The fixed-mantissa factorization \eqref{eq:exact-factorization} was also evaluated
at $y=4/3,6/5,8/7$ and $N=0,1,3,10$ with eighty decimal digits. The largest relative
discrepancy was below $3\times10^{-78}$.

\subsection{Reproducible symbolic code}

The following compact SymPy program implements the resultant route. It is included
as mathematical pseudocode rather than as a software dependency of the document.

\begin{lstlisting}[style=pseudo,language=Python,caption={Orbit-polynomial computation.}]
import sympy as sp
from math import gcd

def order_mod_2(q: int) -> int:
    x = 1
    for d in range(1, 2*q + 1):
        x = (2*x) % q
        if x == 1:
            return d
    raise ArithmeticError("order not found")

def orbit_polynomial(q: int):
    if q <= 1 or q % 2 == 0:
        raise ValueError("q must be odd and > 1")
    z, X = sp.symbols("z X")
    d = order_mod_2(q)
    F = sp.prod(1 - z**(2**j) for j in range(d))
    Phi = sp.cyclotomic_poly(q, z)
    R = sp.Poly(sp.resultant(Phi, X - F, z), X)
    factors = sp.factor_list(R.as_expr(), X)[1]
    P = sp.Integer(1)
    for factor, exponent in factors:
        if exponent % d:
            raise ArithmeticError("resultant is not a d-th power")
        P *= factor**(exponent // d)
    return d, sp.Poly(sp.expand(P), X)

for q in [3, 5, 7, 9, 11, 15, 17, 21, 23, 31]:
    d, P = orbit_polynomial(q)
    print(q, d, sp.factor(P.as_expr()))
\end{lstlisting}

\section{A general-base extension}\label{sec:general-base}

The arithmetic mechanism does not depend on base two until the sharp endpoint
argument. For an integer $b\ge2$, define
\begin{equation}
 \mathscr U_b(x):=\prod_{n=0}^{\infty}
 \sinc\!\left(\frac{\pi x}{b^n}\right).
 \label{eq:Ub}
\end{equation}
For $1<y<2$ and $t=y-1$, put
\begin{equation}
 H_b(y):=\prod_{r=1}^{\infty}\sinc\!\left(\frac{\pi y}{b^r}\right),
 \qquad
 Q_{b,n}(t):=\prod_{j=0}^{n-1}|2\sin(\pi b^jt)|.
 \label{eq:Hb-Qb}
\end{equation}
(The interval $1<y<2$ is enough for a canonical family of rays; one may use
$1<y<b$ with only notational changes.)

\begin{theorem}[Exact rational rays in base $b$]\label{thm:base-b}
For every $N\ge0$,
\begin{equation}
 |\mathscr U_b(b^Ny)|
 =b^{-N(N+1)/2}(\pi y)^{-(N+1)}H_b(y)2^{-(N+1)}Q_{b,N+1}(t).
 \label{eq:base-b-factorization}
\end{equation}
Let $q>1$, $(a,q)=1$, $(b,q)=1$, $d=\ord_q(b)$, and
\begin{equation}
 A_{a,q}^{(b)}:=\left|\prod_{j=0}^{d-1}
 (1-\zeta_q^{ab^j})\right|.
 \label{eq:A-base-b}
\end{equation}
On the ray $x_N=b^N(1+a/q)$ there is an exact periodic law
\begin{equation}
 \log|\mathscr U_b(x_N)|
 =-\frac{(\log x_N)^2}{2\log b}
 -\kappa_{a,q}^{(b)}\log x_N+\Omega_{a,q}^{(b)}(N),
 \label{eq:base-b-ray}
\end{equation}
where $\Omega_{a,q}^{(b)}$ has period dividing $d$ and
\begin{equation}
 \boxed{
 \kappa_{a,q}^{(b)}=\frac12+\frac{\log(2\pi)}{\log b}
 -\frac{\log A_{a,q}^{(b)}}{d\log b}.}
 \label{eq:kappa-base-b}
\end{equation}
If $H_{q,b}=\langle b\rangle\subseteq(\Z/q\Z)^\times$, then the mean over its
cosets is
\begin{equation}
 \boxed{
 \frac1{|G_q/H_{q,b}|}\sum_C\kappa_C^{(b)}
 =\frac12+\frac{\log(2\pi)}{\log b}
 -\frac{\log\Phi_q(1)}{\varphi(q)\log b}.}
 \label{eq:mean-base-b}
\end{equation}

\smallskip\noindent\newresult
\end{theorem}

\begin{proof}
Split the product at $n=N$ exactly as in \cref{thm:exact-factorization}; the
denominator now contributes $b^{N(N+1)/2}$, while the factor $2^{-(N+1)}$ still
comes from replacing $|\sin|$ by $|2\sin|$. If $t=a/q$, multiplication by $b$ has
period $d$ modulo $q$, so $Q_{b,n+d}=A_{a,q}^{(b)}Q_{b,n}$. Completing the square
in $\log x_N=N\log b+\log y$ gives \eqref{eq:kappa-base-b}. The mean law follows
from the same norm identity because the $H_{q,b}$-cosets partition $G_q$.
\end{proof}

\begin{remark}[What is genuinely binary]
The resultant, relative norm, exact periodic remainder, and mean-exponent law are
base-independent. The sharp $\sqrt3$ endpoint uses the special identity
$\cos(2x)=2\cos^2x-1$ and the fact that squaring permutes the cycle. Finding the
corresponding sharp constant for $b\ge3$ is a distinct ergodic-optimization problem.
\end{remark}

\section{A Lean formalization blueprint}\label{sec:lean-blueprint}

The new results are deliberately organized so that most proof obligations are
finite algebra, elementary real analysis, or existing cyclotomic infrastructure.
This section gives a concrete path from the report to the repository's formalized
core.

\subsection{Suggested module boundary}

\begin{description}[style=nextline,leftmargin=2.8em]
 \item[\path{Analysis/FabiusFunction/ThueMorse/CyclotomicOrbit.lean}]
 Define $T_n$, prove its product formula, evaluate it on a root of unity, and prove
 the block recurrence $T_{n+d}=\Gamma T_n$.
 \item[\path{Analysis/FabiusFunction/Fourier/RationalRay.lean}]
 Split the infinite sinc product at $N$, prove \eqref{eq:exact-factorization}, and
 derive the exact periodic log law.
 \item[\path{Analysis/FabiusFunction/NumberTheory/OrbitPolynomial.lean}]
 Define $H_q=\langle2\rangle$, the relative norm, and the orbit polynomial. Prove
 integrality, the norm product, and the mean exponent identity.
 \item[\path{Analysis/FabiusFunction/Analysis/MersenneRay.lean}]
 Establish the exact products \eqref{eq:Aminus-product}--\eqref{eq:Aplus-product},
 the convergence to $h$, and then the asymptotic estimates.
 \item[\path{Analysis/FabiusFunction/QCalculus/RayBridge.lean}]
 Combine the Mersenne factorial identity with the sinc-tail limit to formalize
 \eqref{eq:qbridge-limits}.
\end{description}

\subsection{Core theorem signatures}

The exact identities can be formalized before any logarithms, avoiding extended
real values at zeros. Indicative signatures are:

\begin{lstlisting}[style=pseudo,caption={Indicative Lean theorem interfaces.}]
-- Finite Thue-Morse polynomial at a root of unity.
theorem tmPoly_eval_cycle
    (hq : Odd q) (ha : Nat.Coprime a q) :
    tmPoly (n + orderOfTwo q) (zeta q ^ a) =
      cycleMultiplier a q * tmPoly n (zeta q ^ a)

-- Exact absolute factorization of the Rvachev Fourier product.
theorem upFourier_abs_dyadic_mantissa
    (hy : 1 < y) (hy2 : y < 2) :
    abs (upFourier ((2 : R) ^ N * y)) =
      2 ^ (-(N * (N + 1) / 2 : Z)) *
      (Real.pi * y) ^ (-(N + 1 : Z)) *
      sincTail y * 2 ^ (-(N + 1 : Z)) *
      orbitSineProduct (N + 1) (y - 1)

-- Product over G/H and the cyclotomic value at one.
theorem prod_cycleMultiplier_cosets :
    Finset.prod (allCosets q) (fun C =>
      cycleMultiplier C.rep q) = cyclotomic q 1

-- Sharp periodic-orbit bound.
theorem cycleAmplitude_geomMean_le_sqrt_three :
    Real.rpow (cycleAmplitude a q) (1 / orderOfTwo q) <= Real.sqrt 3
\end{lstlisting}

The displayed syntax is schematic. The main implementation decision is whether to
represent roots of unity in $\C$ or in an abstract cyclotomic extension. The
resultant theorem is easiest abstractly; the absolute-value ray law is easiest in
$\C$ with the standard embedding.

\subsection{Proof dependency order}

A low-risk formalization sequence is:
\begin{enumerate}
 \item finite product and block recurrence;
 \item exact fixed-mantissa factorization without logarithms;
 \item rational preperiod/period dichotomy;
 \item norm product and mean exponent;
 \item sharp $\sqrt3$ inequality, which uses only finite real sums and AM--GM;
 \item analytic Mersenne estimates;
 \item resultant and Ramanujan trace formulas, which can be added independently.
\end{enumerate}
The first five steps already formalize the central qualitative discoveries and do
not depend on delicate asymptotic libraries.

\section{Further frontier problems}\label{sec:open}

The exact arithmetic-ray theory raises a set of focused questions that are sharper
than the generic request to ``understand Fourier decay.''

\subsection{The closure of the ray-exponent spectrum}

\begin{conjecture}[Interval closure]
Let
\[
 \mathcal K_{\mathrm{rat}}=
 \left\{\kappa_{a,q}:q>1\text{ odd},\ (a,q)=1\right\}.
\]
Its closure may be the whole interval $[\kappa_{\min},\infty)$.

\smallskip\noindent\openstatus
\end{conjecture}

Periodic orbit measures are dense among invariant measures of the doubling map for
continuous observables, but $\log|2\sin(\pi t)|$ has logarithmic singularities. The
conjecture therefore requires a controlled specification argument that keeps
periodic words away from zero while prescribing their cycle average. The Mersenne
family gives the unbounded direction, and the denominator-three orbit gives the
lower endpoint; the missing issue is interpolation.

\subsection{Extremizers at fixed period}

For each $d$, determine
\[
 A_d^{\max}:=\max_{\ord_q(2)=d}A_{a,q},
 \qquad
 A_d^{\min}:=\min_{\ord_q(2)=d}A_{a,q}.
\]
The sharp global maximum of $A^{1/d}$ is known from \cref{thm:sqrt3-bound}, but the
finite-period problem may have a rich symbolic structure. A successful solution
would give the sharpest and slowest rational rays at every scale period.

\subsection{Stabilizers and repeated orbit-polynomial roots}

Characterize the subgroup
\[
 \operatorname{Stab}(\Gamma_{1,q})
 =\{u\in G_q:\Gamma_{u,q}=\Gamma_{1,q}\}.
\]
It contains $H_q$ but can be larger, as $q=15$ and $21$ show. Equivalently, determine
when $\mathscr P_{q,2}$ is squarefree and when it is a proper power. This is a
concrete problem in cyclotomic units and automatic polynomial values.

\subsection{Regulator statistics}

When $q$ has at least two prime divisors, \eqref{eq:product-A} says
$\sum_C\log A_C=0$. The vector $(\log A_C)_C$ lies in the logarithmic embedding of
a cyclotomic-unit group. Its covariance, extreme coordinates, and limiting
distribution as $q\to\infty$ should connect ray decay with regulators and unit
lattices. Prime-power denominators add the deterministic shift
$\log p/\varphi(p^s)$ from \eqref{eq:mean-kappa}.

\subsection{Beyond the first Mersenne correction}

The proof of \cref{thm:pm-product-asymptotic} can be iterated. Taylor expansion in
$\delta=2^{-d}$ produces an all-orders transseries whose coefficients are convergent
cotangent/cosecant sums over dyadic angles, together with the tail expansion of
$H_d(1)$. Determining a closed algebra for these coefficients would parallel the
periodic saddle-jet machinery already developed for the small-argument Fabius
asymptotic.

\subsection{A wider $q$-binomial bridge}

The bridge theorem uses the terminal Mersenne factor $2^d-1$ and the full product
$\mathfrak M_d=\prod_{j\le d}(2^j-1)$. A broader goal is to connect general
cyclotomic orbit norms $A_{a,q}$ to coefficient extraction in the half-base
$q$-binomial formulas. Possible entry points are:
\begin{itemize}
 \item denominators dividing products of Mersenne integers;
 \item discrete Fourier transforms of the Gaussian interpolation blocks at roots
 of unity;
 \item resultant identities in which $F_d(z)$ is paired with the half-base
 Pochhammer polynomial.
\end{itemize}
A direct finite identity here would join the exact dyadic values of $F$ to the
arithmetic rays of $\widehat{\Up}$, rather than linking them only asymptotically.

\subsection{Sharp constants in other bases}

For $b\ge3$, determine
\[
 \sup_{q,a}\left(A_{a,q}^{(b)}\right)^{1/\ord_q(b)}
\]
and classify the maximizing periodic orbits. The binary proof is exceptional
because the real part of $z\mapsto z^2$ closes quadratically. The general problem
looks like a finite-dimensional ergodic-optimization problem for
$\log|1-z|$ under $z\mapsto z^b$.

\section{Conclusion}\label{sec:conclusion}

The reviewed corpus already contains two powerful but largely parallel bodies of
work: exact $q=1/2$ arithmetic for Fabius values and dynamical Fourier decay for the
Rvachev product. The arithmetic-ray viewpoint inserts the finite Thue--Morse
polynomial between them. On rational mantissas, the doubling cocycle closes; its
cycle multiplier is a relative cyclotomic norm, and the Fourier decay becomes an
exact log-quadratic formula with a finite periodic remainder.

This perspective yields several rigid structures at once:
\begin{itemize}
 \item an integral orbit polynomial and a perfect-power resultant;
 \item Ramanujan-sum formulas for every multiplier trace;
 \item a one-integer mean-exponent law governed by $\Phi_q(1)$;
 \item a sharp rational-ray minimum from the orbit $\{1/3,2/3\}$;
 \item unbounded exponents from Mersenne denominators;
 \item a twin Mersenne/plus-denominator expansion with opposite first corrections;
 \item a direct limit identity joining cyclotomic norms, the half-base
 $q$-Pochhammer constant, and $-\mathscr U'(1)$.
\end{itemize}
The results are elementary enough to be formalized in stages, yet arithmetic enough
to open new questions about cyclotomic units, regulators, periodic-orbit spectra,
and Gaussian $q$-binomial transforms. They therefore fit the repository's intended
frontier role: mathematically testable statements, clearly separated from imported
results, and organized for eventual migration into Lean.

\appendix

\section{Corpus inventory}\label{app:inventory}

The audit counted fifty-seven TeX files: twenty-five in the living documentation
tree and thirty-two frozen historical versions. The list below records the path
relative to \path{Analysis/FabiusFunction/docs}. Historical items are grouped by
the archive directory; filenames follow the base name shown.

\subsection{Living documentation: 25 TeX files}

\begin{longtable}{@{}r p{0.88\textwidth}@{}}
\toprule
No. & Relative path \\
\midrule
\endfirsthead
\toprule
No. & Relative path \\
\midrule
\endhead
1 & \path{FabiusFunction_Mathematical_Glossary/FabiusFunction_Mathematical_Glossary.tex} \\
2 & \path{Fabius_Function_and_Rvachev_Up/Fabius_Function_and_Rvachev_Up.tex} \\
3 & \path{Non_Elementarity_of_the_Fabius_Function/Non_Elementarity_of_the_Fabius_Function.tex} \\
4 & \path{fabius_lean_walkthrough/fabius_lean_walkthrough.tex} \\
5 & \path{non-formalized-research-frontiers/Non_Formalized_Research_Frontiers.tex} \\
6 & \path{non-formalized-research-frontiers/drafts/Fabius_Thue_Morse_Formula_Atlas/Fabius_Thue_Morse_Formula_Atlas.tex} \\
7 & \path{non-formalized-research-frontiers/drafts/rvachev_up_fourier_decay/Rvachev_Up_Fourier_Decay-2/Rvachev_Up_Fourier_Decay.tex} \\
8 & \path{non-formalized-research-frontiers/drafts/rvachev_up_fourier_decay/Rvachev_Up_Fourier_Decay-3/Fourier_decay_of_Rvachev_up.tex} \\
9 & \path{non-formalized-research-frontiers/drafts/rvachev_up_fourier_decay/Rvachev_Up_Fourier_Decay_Comparative_Audit/Rvachev_Up_Fourier_Decay_Comparative_Audit.tex} \\
10 & \path{non-formalized-research-frontiers/drafts/rvachev_up_fourier_decay/Rvachev_Up_Fourier_Decay_Gentle_Guide/Rvachev_Up_Fourier_Decay_Gentle_Guide.tex} \\
11 & \path{non-formalized-research-frontiers/drafts/rvachev_up_fourier_decay/Rvachev_Up_Fourier_Decay_Second_Wave_Audit/Rvachev_Up_Fourier_Decay_Second_Wave_Audit.tex} \\
12 & \path{non-formalized-research-frontiers/drafts/rvachev_up_fourier_decay/asymptotic-decay-rate-of-an-infinite-product-of-sinc-functions/asymptotic-decay-rate-of-an-infinite-product-of-sinc-functions.tex} \\
13 & \path{non-formalized-research-frontiers/drafts/rvachev_up_fourier_decay/rvachev_up_fourier_decay-1/rvachev_up_fourier_decay.tex} \\
14 & \path{non-formalized-research-frontiers/drafts/rvachev_up_fourier_decay/rvachev_up_fourier_decay-4/rvachev_up_fourier_decay.tex} \\
15 & \path{non-formalized-research-frontiers/drafts/rvachev_up_fourier_decay/rvachev_up_fourier_decay-5/Rvachev_Up_Fourier_Decay_Final.tex} \\
16 & \path{non-formalized-research-frontiers/drafts/rvachev_up_fourier_decay/rvachev_up_fourier_decay-6/Rvachev_Up_Fourier_Decay_Final_Synthesis.tex} \\
17 & \path{non-formalized-research-frontiers/drafts/rvachev_up_fourier_decay/rvachev_up_fourier_decay-7/rvachev_up_fourier_decay_final.tex} \\
18 & \path{non-formalized-research-frontiers/drafts/rvachev_up_fourier_decay/rvachev_up_fourier_decay-8/rvachev_up_fourier_decay_final.tex} \\
19 & \path{papers/AIF_2017__67_2_637_0/AIF_2017__67_2_637_0.tex} \\
20 & \path{papers/AIF_2017__67_2_637_0/AIF_2017__67_2_637_0-corrected.tex} \\
21 & \path{papers/arXiv-1502.06738v2/Lacunary_products.tex} \\
22 & \path{papers/arXiv-1702.05442v1/09-Function.tex} \\
23 & \path{papers/arXiv-1702.06487v3/157-Arithmetic-v3.tex} \\
24 & \path{papers/arXiv-2211.13570v1/document.tex} \\
25 & \path{papers/ubiq15/ubiq15-corrected.tex} \\
\bottomrule
\end{longtable}

\subsection{Frozen archive: 32 TeX files}

\begin{longtable}{@{}r p{0.88\textwidth}@{}}
\toprule
No. & Relative path under \path{archive/} \\
\midrule
\endfirsthead
\toprule
No. & Relative path under \path{archive/} \\
\midrule
\endhead
1 & \path{early-notes/Fabius Asymptotic/Fabius Asymptotic.tex} \\
2 & \path{early-notes/K-fold summation over the signed Thue-Morse sequence/K-fold summation over the signed Thue-Morse sequence.tex} \\
3 & \path{standalone-studies/Non_Elementarity_of_the_Fabius_Function/Non_Elementarity_of_the_Fabius_Function.tex} \\
4 & \path{standalone-studies/Small_Argument_Asymptotics/Small_Argument_Asymptotics.tex} \\
5 & \path{standalone-studies/fabius-inverse-dyadic-closed-form/fabius-inverse-dyadic-closed-form.tex} \\
6 & \path{primary-exposition/Fabius_Function_and_Rvachev_Up/Fabius_Function_and_Rvachev_Up.tex} \\
7 & \path{primary-exposition/Fabius_Function_and_Rvachev_Up-2/Fabius_Function_and_Rvachev_Up.tex} \\
8 & \path{primary-exposition/Fabius_Function_and_Rvachev_Up-3/Fabius_Function_and_Rvachev_Up.tex} \\
9 & \path{primary-exposition/Fabius_and_Rvachev_Up-2/Fabius_and_Rvachev_Up.tex} \\
10 & \path{primary-exposition/Fabius_and_Rvachev_up/Fabius_and_Rvachev_up.tex} \\
11 & \path{primary-exposition/Fabius_function_and_up_function/Fabius_function_and_up_function.tex} \\
12 & \path{primary-supplements/Fabius_Function_Missing_Parts/Fabius_Function_Missing_Parts.tex} \\
13 & \path{primary-supplements/Fabius_Function_Missing_Parts-2/Fabius_Function_Missing_Parts.tex} \\
14 & \path{q-calculus/Fabius_q_Special_Functions/Fabius_q_Special_Functions.tex} \\
15 & \path{q-calculus/fabius-q-special-functions/fabius-q-special-functions.tex} \\
16 & \path{lean-walkthrough/Fabius_Function_Lean_Walkthrough/Fabius_Function_Lean_Walkthrough.tex} \\
17 & \path{lean-walkthrough/fabius_lean_walkthrough/fabius_lean_walkthrough.tex} \\
18 & \path{glossary/FabiusFunction_Mathematical_Glossary/FabiusFunction_Mathematical_Glossary.tex} \\
19 & \path{glossary/FabiusFunction_Mathematical_Glossary-2/FabiusFunction_Mathematical_Glossary.tex} \\
20 & \path{research-frontiers/Fabius_Dyadic_Asymptotic_Bridge/Fabius_Dyadic_Asymptotic_Bridge.tex} \\
21 & \path{research-frontiers/Fabius_Dyadic_Formulae_and_Alternative_Representations/Fabius_Dyadic_Formulae_and_Alternative_Representations.tex} \\
22 & \path{research-frontiers/Fabius_Dyadic_Formulae_to_Asymptotics/Fabius_Dyadic_Formulae_to_Asymptotics.tex} \\
23 & \path{research-frontiers/Fabius_Dyadic_q_Connections/Fabius_Dyadic_q_Connections.tex} \\
24 & \path{research-frontiers/Fabius_Thue_Morse_Convergence_Rate/Fabius_Thue_Morse_Convergence_Rate.tex} \\
25 & \path{research-frontiers/Fabius_Thue_Morse_Convergence_Rates/Fabius_Thue_Morse_Convergence_Rates.tex} \\
26 & \path{research-frontiers/Primary_Exposition_Gap_Register/Primary_Exposition_Gap_Register.tex} \\
27 & \path{research-frontiers/Repeated_Integration_and_Rvachev_Up/Repeated_Integration_and_Rvachev_Up.tex} \\
28 & \path{research-frontiers/Rvachev_Up_from_Repeated_Integration/Rvachev_Up_from_Repeated_Integration.tex} \\
29 & \path{research-frontiers/non-formalized-research-frontiers/non-formalized-research-frontiers.tex} \\
30 & \path{research-frontiers/Thue_Morse_Formula_Atlas-1/Thue_Morse_Formula_Atlas.tex} \\
31 & \path{research-frontiers/Thue_Morse_Formula_Atlas-2/Consolidated_Formula_Atlas_for_the_Thue_Morse_Sequence.tex} \\
32 & \path{research-frontiers/Thue_Morse_Formula_Atlas-3/Thue_Morse_Formula_Atlas.tex} \\
\bottomrule
\end{longtable}

\section{Constants and numerical reference}\label{app:constants}

For convenient independent checking:
\begin{align*}
 \kappa_0&=3.1514961294723187980432792951080073350\ldots,\\
 \kappa_{\min}&=2.3590148791117407073164098231340990806\ldots,\\
 h&=0.5537712758888100953442174163445488813\ldots,\\
 -\log_2h&=0.8526378718208337174366143795551934080\ldots,\\
 C_{\cot}&=-1.2837124730461260445945438959726042474\ldots,\\
 (1/2;1/2)_\infty&=0.2887880950866024212788997219292307801\ldots,\\
 \mathcal C_*&=0.1599225518776068325779014528557986690\ldots.
\end{align*}

\begin{thebibliography}{99}

\bibitem{proveit}
V.~Reshetnikov,
\emph{ProveIt: Analysis/FabiusFunction}, living Lean development and documentation,
GitHub repository, 2026.
\url{https://github.com/VladimirReshetnikov/ProveIt/tree/main/Analysis/FabiusFunction}.

\bibitem{ariascompact}
J.~Arias de Reyna,
\emph{An infinitely differentiable function with compact support: Definition and
properties}, English translation, arXiv:1702.05442, 2017.

\bibitem{ariasarith}
J.~Arias de Reyna,
\emph{Arithmetic of the Fabius function}, Integers \textbf{18} (2018), Paper A51;
arXiv:1702.06487.

\bibitem{aistleitner2017}
C.~Aistleitner, R.~Hofer, and G.~Larcher,
\emph{On evil Kronecker sequences and lacunary trigonometric products},
Annales de l'Institut Fourier \textbf{67} (2017), no.~2, 637--687,
\url{https://doi.org/10.5802/aif.3094}.

\bibitem{alloucheshallit1999}
J.-P.~Allouche and J.~Shallit,
\emph{The ubiquitous Prouhet--Thue--Morse sequence}, in
\emph{Sequences and Their Applications}, Springer Series in Discrete Mathematics
and Theoretical Computer Science, 1999, 1--16.

\bibitem{alloucheshallitbook}
J.-P.~Allouche and J.~Shallit,
\emph{Automatic Sequences: Theory, Applications, Generalizations},
Cambridge University Press, 2003.

\bibitem{sobolewski2019}
B.~Sobolewski,
\emph{Cyclotomic properties of polynomials associated with automatic sequences},
arXiv:1909.12052, 2019.

\bibitem{toth2022}
L.~T\'oth,
\emph{Linear combinations of Dirichlet series associated with the Thue--Morse
sequence}, arXiv:2211.13570, 2022.

\bibitem{hardywright}
G.~H.~Hardy and E.~M.~Wright,
\emph{An Introduction to the Theory of Numbers}, 6th ed., Oxford University Press,
2008.

\bibitem{washington}
L.~C.~Washington,
\emph{Introduction to Cyclotomic Fields}, 2nd ed., Graduate Texts in Mathematics
83, Springer, 1997.

\end{thebibliography}

\end{document}
